\documentclass[11pt,a4paper]{article}

% Pakete
\usepackage[utf8]{inputenc}
\usepackage[T1]{fontenc}
\usepackage[ngerman]{babel}
\usepackage[left=2cm,right=2cm,top=2cm,bottom=2cm]{geometry}
\usepackage{helvet}
\renewcommand{\familydefault}{\sfdefault}
\usepackage{amsmath,amssymb}
\usepackage{graphicx}
\usepackage{tikz}
\usetikzlibrary{arrows.meta,positioning,decorations.pathmorphing}
\usepackage{tcolorbox}
\tcbuselibrary{skins,breakable}
\usepackage{enumitem}
\usepackage{tabularx}
\usepackage{colortbl}
\usepackage{siunitx}
\sisetup{locale=DE, per-mode=fraction}

% Fachfarbe Physik
\definecolor{physik}{HTML}{2c5aa0}
\definecolor{physiklight}{HTML}{e3f2fd}
\definecolor{warnung}{HTML}{b35c00}
\definecolor{warnunglight}{HTML}{fff3e0}

% Box-Stile
\tcbset{
    infobox/.style={
        colback=physiklight,
        colframe=physik,
        fonttitle=\bfseries,
        title=#1,
        boxrule=1pt,
        arc=2pt,
        left=5pt,
        right=5pt,
        top=5pt,
        bottom=5pt
    },
    merksatz/.style={
        colback=physiklight,
        colframe=physik,
        boxrule=2pt,
        arc=2pt,
        left=8pt,
        right=8pt,
        top=8pt,
        bottom=8pt
    },
    arbeitsblatt/.style={
        colback=warnunglight,
        colframe=warnung,
        fonttitle=\bfseries,
        title=Arbeitsblatt,
        boxrule=1pt,
        arc=2pt,
        left=5pt,
        right=5pt,
        top=5pt,
        bottom=5pt
    }
}

% Kopfzeile
\usepackage{fancyhdr}
\pagestyle{fancy}
\fancyhf{}
\lhead{\small Physik 12 GK}
\chead{\small Lernblatt}
\rhead{\small 09.01.2026}
\renewcommand{\headrulewidth}{0.4pt}

% Keine Einrückung
\setlength{\parindent}{0pt}
\setlength{\parskip}{6pt}

\begin{document}

% Titel
\begin{center}
{\Large\textbf{Lernblatt: Gegenfeld und Lichtquanten}}\\[0.3em]
{\normalsize Stunde 2 | Freitag, 09.01.2026 | Doppelstunde}
\end{center}

\vspace{0.3em}

% ============================================================
% WIEDERHOLUNG
% ============================================================

\section*{1. Wiederholung: Was wissen wir?}

Letzte Stunde haben wir den \textbf{Hallwachs-Versuch} kennengelernt:

\begin{itemize}[nosep]
\item UV-Licht auf negativ geladene Zinkplatte $\rightarrow$ Ladung nimmt ab (Elektronen werden herausgelöst)
\item Sichtbares Licht $\rightarrow$ \textbf{keine Wirkung}, egal wie intensiv!
\end{itemize}

\begin{tcolorbox}[infobox={Problem: Widerspruch zum Wellenmodell}]
\textbf{Wellenmodell sagt:} Höhere Intensität = mehr Energie $\rightarrow$ müsste Elektronen auslösen.

\textbf{Experiment zeigt:} Intensität hilft nicht! Nur die \textbf{Frequenz} (Farbe) entscheidet.

$\Rightarrow$ Das Wellenmodell kann den Photoeffekt \textbf{nicht erklären}!
\end{tcolorbox}

% ============================================================
% GEGENFELDMETHODE
% ============================================================

\section*{2. Die Gegenfeldmethode}

\textbf{Frage:} Wie schnell sind die herausgelösten Elektronen?

\textbf{Idee:} Wir bremsen die Elektronen mit einem elektrischen Feld ab.

\begin{center}
\begin{tikzpicture}[scale=0.9]
    % Fotozelle
    \draw[thick] (0,0) rectangle (6,4);
    \node at (3,4.3) {\small Vakuumfotozelle};

    % Kathode
    \draw[very thick, physik] (1,0.5) -- (1,3.5);
    \node[left] at (0.8,2) {\small Kathode};
    \node[left] at (0.8,1.5) {\small (Metall)};

    % Anode
    \draw[very thick] (5,0.5) -- (5,3.5);
    \node[right] at (5.2,2) {\small Anode};

    % Licht
    \draw[->, thick, yellow!80!orange, decorate, decoration={snake, amplitude=2pt, segment length=8pt}] (-1.5,2.5) -- (0.8,2.5);
    \node[left] at (-1.5,2.5) {\small Licht ($f$)};

    % Elektronen
    \draw[->, thick, blue] (1.3,2.5) -- (2.5,2.5);
    \draw[->, thick, blue] (1.3,1.5) -- (2.2,1.5);
    \node[blue] at (2,3) {\small $e^-$};

    % Gegenfeld
    \draw[<-, thick, red] (3,1) -- (4.5,1);
    \node[red] at (3.7,0.6) {\small E-Feld bremst};

    % Spannungsquelle
    \draw (1,-0.5) -- (1,-1.5) -- (2.5,-1.5);
    \draw (5,-0.5) -- (5,-1.5) -- (3.5,-1.5);
    \draw (2.5,-1.8) rectangle (3.5,-1.2);
    \node at (3,-1.5) {\small $U$};
    \node at (2.2,-1.5) {\small $-$};
    \node at (3.8,-1.5) {\small $+$};

    % Amperemeter
    \draw (5,-0.5) -- (7,-0.5) -- (7,2);
    \draw (7,2) circle (0.4);
    \node at (7,2) {\small A};
    \draw (7,2.4) -- (7,4.5) -- (1,4.5) -- (1,3.7);
\end{tikzpicture}
\end{center}

\textbf{Prinzip:}
\begin{enumerate}[nosep]
\item Licht löst Elektronen aus der Kathode
\item Elektronen fliegen zur Anode $\rightarrow$ Strom fließt
\item Wir erhöhen die \textbf{Gegenspannung} $U$
\item Bei der \textbf{Grenzspannung} $U_G$: Strom wird Null (alle Elektronen gestoppt)
\end{enumerate}

\begin{tcolorbox}[merksatz]
\textbf{Energiebilanz:} Die kinetische Energie der schnellsten Elektronen wird vollständig in elektrische Energie umgewandelt:

\vspace{0.3em}
\hfill $\boxed{E_{\text{kin,max}} = e \cdot U_G}$ \hfill\mbox{}
\vspace{0.3em}

\textbf{Einheit:} Wenn $U_G$ in Volt gemessen wird, ist $E_{\text{kin}}$ numerisch gleich in \textbf{Elektronenvolt (eV)}.
\end{tcolorbox}

\begin{tcolorbox}[arbeitsblatt]
$\rightarrow$ \textbf{Jetzt:} Öffne die PhET-Simulation und nimm eine Messreihe auf (AB-02, Aufgabe 2).
\end{tcolorbox}

\newpage

% ============================================================
% MESSERGEBNISSE
% ============================================================

\section*{3. Was zeigen die Messergebnisse?}

Nach der Messreihe in der Simulation erkennst du:

\begin{tcolorbox}[infobox={Zentrale Beobachtungen}]
\begin{enumerate}
\item Mit zunehmender \textbf{Beleuchtungsstärke} nimmt die \textbf{Anzahl} der Elektronen zu (mehr Strom), aber \textbf{nicht ihre Energie}!

\item Die \textbf{kinetische Energie} der Elektronen hängt \textbf{nur von der Frequenz} des Lichts ab.

\item Die Energie wächst \textbf{linear} mit der Frequenz: $E_{\text{kin}} \sim f$
\end{enumerate}
\end{tcolorbox}

% ============================================================
% EINSTEINS ERKLÄRUNG
% ============================================================

\section*{4. Einsteins Lichtquanten-Hypothese (1905)}

Albert Einstein fand die Erklärung:

\begin{tcolorbox}[merksatz]
\textbf{Licht besteht aus einzelnen Energiepaketen} -- den \textbf{Lichtquanten} oder \textbf{Photonen}.

Jedes Photon trägt die Energie:

\vspace{0.5em}
\hfill $\boxed{E_{\text{Photon}} = h \cdot f}$ \hfill\mbox{}
\vspace{0.5em}

\begin{tabular}{ll}
$h$ & = \SI{6,63e-34}{J.s} = \SI{4,14e-15}{eV.s} \quad (Plancksches Wirkungsquantum) \\
$f$ & = Frequenz des Lichts in Hz \\
\end{tabular}
\end{tcolorbox}

\begin{tcolorbox}[infobox={Kernidee: Energie ist gequantelt!}]
\begin{itemize}[nosep]
\item Ein Photon = \textbf{ein} Energiepaket = $h \cdot f$
\item \textbf{Unteilbar} -- ganz oder gar nicht (wie eine Münze)
\item Mehr Intensität = mehr Photonen, \textbf{nicht} mehr Energie pro Photon
\end{itemize}
\end{tcolorbox}

\textbf{Warum erklärt das den Photoeffekt?}

\begin{itemize}[nosep]
\item Ein Photon gibt seine \textbf{gesamte Energie} an \textbf{ein} Elektron ab
\item Mehr Intensität = mehr Photonen, aber \textbf{nicht} mehr Energie pro Photon
\item Höhere Frequenz = mehr Energie pro Photon $\rightarrow$ schnellere Elektronen
\end{itemize}

% ============================================================
% PHOTOGLEICHUNG
% ============================================================

\section*{5. Die Photogleichung}

Ein Teil der Photonenenergie wird zum Herauslösen des Elektrons benötigt:

\begin{tcolorbox}[merksatz]
\textbf{Einsteinsche Gleichung (Photogleichung):}

\vspace{0.5em}
\hfill $\boxed{h \cdot f = E_{\text{kin,max}} + W_A}$ \hfill oder \hfill $\boxed{E_{\text{kin,max}} = h \cdot f - W_A}$ \hfill\mbox{}
\vspace{0.5em}

\begin{tabular}{ll}
$W_A$ & = \textbf{Austrittsarbeit} (materialabhängig) \\
& = Mindestenergie zum Herauslösen eines Elektrons \\
\end{tabular}
\end{tcolorbox}

\textbf{Grenzfrequenz:} Bei $f = f_G$ ist gerade $E_{\text{kin}} = 0$:

\hfill $h \cdot f_G = W_A$ \hfill $\Rightarrow$ \hfill $f_G = \dfrac{W_A}{h}$ \hfill\mbox{}

\vspace{0.5em}

\begin{tcolorbox}[arbeitsblatt]
$\rightarrow$ \textbf{Jetzt:} Nimm die Messreihe auf und zeichne die Einsteinsche Gerade (AB-02, Aufgaben 1--3).
\end{tcolorbox}

% ============================================================
% EINSTEIN-GERADE
% ============================================================

\section*{6. Die Einsteinsche Gerade}

Wenn du deine Messwerte in ein Diagramm ($E_{\text{kin}}$ gegen $f$) einträgst, erkennst du:

\begin{center}
\begin{tikzpicture}[scale=0.9]
    % Diagramm mit Extrapolation zu -W_A
    % Natrium: W_A = 2,3 eV, f_G = 5,56·10^14 Hz

    % Achsen (y geht bis -W_A)
    \draw[->] (0,0) -- (8.5,0) node[right] {$f$ in $10^{14}$ Hz};
    \draw[->] (0,-2.8) -- (0,1.8) node[above] {$E_{\text{kin}}$ in eV};

    % Achsenbeschriftung
    \foreach \x in {2,4,6,8} {
        \draw (\x,0.08) -- (\x,-0.08) node[below] {\small \x};
    }
    \draw (0.08,1) -- (-0.08,1) node[left] {\small 1};
    \draw (0.08,-1) -- (-0.08,-1) node[left] {\small $-1$};
    \draw (0.08,-2) -- (-0.08,-2) node[left] {\small $-2$};
    \node[below left] at (0,0) {\small 0};

    % Extrapolation (gestrichelt) von (0,-2.3) bis f_G
    \draw[dashed, thick, gray] (0,-2.3) -- (5.56,0);
    % Messgerade (durchgezogen) ab f_G
    \draw[very thick, physik] (5.56,0) -- (8,1.01);

    % Messpunkte
    \fill[physik] (6,0.18) circle (2.5pt);
    \fill[physik] (7,0.60) circle (2.5pt);
    \fill[physik] (8,1.01) circle (2.5pt);

    % Grenzfrequenz f_G
    \fill[red] (5.56,0) circle (3pt);
    \node[below, red] at (5.56,-0.2) {\small $f_G$};

    % y-Achsenabschnitt -W_A
    \fill[orange] (0,-2.3) circle (3pt);
    \node[right, orange] at (0.15,-2.3) {\small $-W_A$};

    % Steigungsdreieck
    \draw[thin] (6,0.18) -- (7,0.18) -- (7,0.60);
    \node[below] at (6.5,0.18) {\scriptsize $\Delta f$};
    \node[right] at (7,0.39) {\scriptsize $\Delta E$};
    \node[right, physik] at (7.2,0.85) {\small Steigung $= h$};
\end{tikzpicture}
\end{center}

\begin{tcolorbox}[merksatz]
\textbf{Interpretation der Einsteinschen Geraden:}

\vspace{0.3em}
\begin{tabular}{ll}
\textbf{Steigung} & $= h$ (Plancksches Wirkungsquantum) \\
\textbf{x-Achsenabschnitt} & $= f_G$ (Grenzfrequenz) \\
\textbf{y-Achsenabschnitt} & $= -W_A$ (negative Austrittsarbeit) \\
\textbf{Geradengleichung} & $E_{\text{kin}} = h \cdot f - W_A$ (Photogleichung!) \\
\end{tabular}

\vspace{0.3em}
$\Rightarrow$ Aus der Steigung kannst du $h$ experimentell bestimmen!
\end{tcolorbox}

\end{document}
